\batchmode
\documentclass{article}
\usepackage[letterpaper, margin=1in]{geometry}
\begin{document}

\section{A chunk of my lecture notes}

One obvious suspect in the poor results above is the crude, first-order accuracy of the forward Euler method.
 The local truncation error of $\mathcal{O}(\Delta t^2)$ means that the \emph{global} error, accumulated over many steps to reach a fixed time $T = N\Delta t$, will scale as $\mathcal{O}(\Delta t)$.
 For many applications, this is simply not good enough.
A logical next step would be to find a method that accounts for higher-order terms in the Taylor series expansion of $\mathbf{y}(t+\Delta t)$.


Expanding $\mathbf{y}(t+\Delta t)$ to higher order, this time going back to a scalar example to keep the notation clean:
\begin{equation}
    y(t+\Delta t) = y(t) + \Delta t \dot{y}(t) + \frac{\Delta t^2}{2} \ddot{y}(t) + \mathcal{O}(\Delta t^3).
\end{equation}
We know that $\dot{y} = f(y,t)$.
Adopting the notation in which $f_y$ and $f_t$ mean the derivative of $f$ with respect to the subscripted variable, we can find the second derivative, $\ddot{y}$, by applying the chain rule to $f$:
\begin{equation}
    \ddot{y} = \frac{d}{dt} f(y(t), t) = \frac{\partial f}{\partial y}\frac{dy}{dt} + \frac{\partial f}{\partial t} = f_y f + f_t.
\end{equation}
Substituting these into the Taylor expansion gives us a second-order accurate update rule:
\begin{equation}
    y(t+\Delta t) = y(t) + \Delta t f + \frac{\Delta t^2}{2} (f_y f + f_t) + \mathcal{O}(\Delta t^3).
\end{equation}
While this ``Taylor series method'' is indeed more accurate, it's often impractical.
It requires us to analytically calculate the partial derivatives of $\mathbf{f}$, which can be extremely complicated for some of the complex functions that arise in physical simulations.
So, while tempting in its simplicity, such Taylor series methods are very rarely used in practice.

\subsection{An artificial subsection}

One of the most important features of Hamiltonian dynamics is that it is time-reversible -- do the integrators we are using so far have this property?
Let's explicitly show that the answer is ``clearly not.''
To demonstrate, let's return to the forward Euler method in the context of a simple harmonic oscillator in one dimension.
The state vector will just be $\vec{y} = \{q, p\}$, and (choosing a convenient system of units) the time derivative is $\vec{f} = \{p, -q\}$. 
We'll probe the time-(ir)reversibility of the method by (1) applying the forward Euler rule forward by one step from some initial state, and then (2) applying the rule again, but in the backwards ($-\Delta t$) time direction; we'll see if the system ends up where it started.

Starting at $\vec{y}_0$, the forward step is just 
\begin{equation}
    \vec{y}_1 = \begin{pmatrix}
    q_0 + \Delta t p_0\\
    p_0 - \Delta t q_0\\
    \end{pmatrix}.
\end{equation}
The final position after a backwards step is then $\vec{y}_f = \vec{y}_1 +(-\Delta t) f(\vec{y}_1)$.
Evaluating this expression, we find
\begin{equation}
    \vec{y}_f = \begin{pmatrix}
    q_0 + \Delta t p_0 - \Delta t p_0 + \Delta t^2 q_0\\
    p_0 - \Delta t q_0 + \Delta t q_0 + \Delta t^2 p_0\\
    \end{pmatrix}
    = (1+\Delta t^2) \begin{pmatrix}
    q_0\\
    p_0\\
    \end{pmatrix}.
\end{equation}
That is: after running time forward and then backwards, we find our system at a state different from where it started.
Surprisingly, just a tiny adjustment to our forward Euler method can correct this flaw.


In the context of a system governed by the Hamiltonian $\mathcal{H} = T+V$ (i.e., composed of a kinetic term that depends only on particle momenta and a potential term that depends only on positions) the symplectic Euler method is formulated as either
\begin{equation}\label{eq:symEul1}
   \begin{matrix}
   p_{n+1} = \\
   q_{n+1} = \\
   \end{matrix}
   \  \
   \begin{matrix}
   p_n - \Delta t V'(q_n)\\
   \ \ \ q_n + \Delta t T'(p_{n+1})\\
   \end{matrix}
\end{equation}
or
\begin{equation}\label{eq:symEul2}
   \begin{matrix}
       q_{n+1} = \\
       p_{n+1} = \\
   \end{matrix}
   \  \
   \begin{matrix}
       q_n + \Delta t T'(p_{n})\\
   \ \ \ p_n - \Delta t V'(q_{n+1})\\
   \end{matrix}
\end{equation}
Notice that in the first formulation, advancing the position requires knowing the \emph{future} momenta (and vice verse in the second formulation).
Notice, furthermore, that these two formulations do the same operations but in the reverse order: a ``kick and then drift'' or a ``drift and then kick'' of the particles.
Because of the nicely separable structure of these equations, though, implementing this is straightforward:



For any function $\rho(\mathbf{q},\mathbf{p})$ of classical phase space variables, the evolution of that function governed by a Hamiltonian $\mathcal{H}$ is given by the Poisson bracket:
\begin{equation}
    \frac{d\rho}{dt} = \{ \rho,\mathcal{H} \} = \sum_i^N \frac{\partial\rho}{\partial\vec{q}_i}\frac{\partial\mathcal{H}}{\partial \vec{p}_i} - \frac{\partial\rho}{\partial\vec{p}_i}\frac{\partial\mathcal{H}}{\partial\vec{q}_i}
\end{equation}
For the purposes of writing down a formal solution, we define the Liouvillian operator, $L_{\mathcal{H}} \rho = \{ \rho, \mathcal{H} \}$.
This lets us cast the time evolution as a simple equation with the formal solution:
\begin{equation}
    \frac{d\rho}{dt} = L_{\mathcal{H}} \rho \quad \Rightarrow \rho(t) = e^{t L_{\mathcal{H}}} \rho(t=0).
\end{equation}

Restricting our attention to Hamiltonians of the form $\mathcal{H} = T(\mathbf{p}) + V(\mathbf{q})$. The properties of the Poisson bracket mean that we can decompose the Liouvillian:
\begin{equation}
    L_{\mathcal{H}} \rho = (L_T+L_V)\rho = \{ \rho, T \} + \{ \rho, V \}.
\end{equation}
Each of these parts are exactly solvable, albeit only accounting for part of the system dynamics.

Sadly, the evolution operator isn't $L_{\mathcal{H}}$ but $e^{t L_{\mathcal{H}}}$, and the Baker-Campbell-Hausdorff (BCH) formula tells us that we are not allowed to just split the evolution into a part that talks only to the positions and a part that talks only to the momenta.
\end{document}
